\documentclass[11pt]{article}
\usepackage[margin=1in]{geometry}
\usepackage{url}

\title{Low-Token Game Development with Failure-Triggered Planning}
\author{}
\date{}

\begin{document}
\maketitle

\section{Introduction}

Consider \emph{Horror Signal Lost}, one of the game requests evaluated in this report. The requested player must tune radio signals, place pins on a map, and advance the story; a Godot project can start successfully while omitting that interaction loop. An agent building such a game must create the project, connect its scenes and scripts, and check both technical startup and the requested behavior. It may also plan before producing a testable project, inspect its own work repeatedly, and carry an expanding conversation into later model calls. Each extra call produces output tokens and may process the accumulated prompt again. The practical efficiency question is when a planning call is worth its token cost.

Planning-after-Trial (PaT) addresses this question in function-level code generation by testing a direct attempt before invoking a planner \cite{yoon2026pat}. A successful trial avoids planning overhead; failure supplies a reason to invest in decomposition. A Godot game presents a related but broader verification problem: its output is a multi-file project, and successful engine startup does not establish that the requested interactions exist. The Low-Token Game Development (LTGD) Agent System adapts PaT's failure-triggered scheduling principle to this setting. Its objective is to avoid planning that a directly generated game does not need while retaining a path for repairing observed defects.

LTGD keeps the user's request in a Pi conversation and assigns three control roles. The Generator writes a first complete Godot project and hands off its actual directory. At the end of that turn, the Executor checks project structure, runs Godot import, and performs a bounded headless startup. Those engine checks do not call a language model. If they pass, the Executor makes a separate model request to compare the original requirements with the current project files. A Godot failure or a concrete missing requirement triggers the Planner, which receives the latest failure evidence and relevant project context, returns focused repair steps, and hands them back to the Generator. The revised project then goes through the same checks. A first-pass success therefore skips the Planner call, although the requirement review still incurs a model call.

This ordering offers two specific routes to lower token use. The Planner is called only for a confirmed problem, and its request contains selected failure evidence rather than the Generator's full conversation. The Generator is also instructed to stop after a focused first pass and to repair only confirmed defects, which may prevent further speculative turns and repeated context processing. These are design mechanisms, not a guarantee of lower total usage. Requirement review reads project text, and a failed trial can add both planning and generation calls. LTGD also uses the selected model for these roles rather than PaT's heterogeneous model allocation \cite{yoon2026pat}. Its net token cost must therefore be measured alongside task completion.

This report asks whether LTGD records lower model usage than ordinary Pi generation for the same game request, and what evidence is available that the resulting project works. It compares model calls, uncached input, cached input reads, output, provider-reported total tokens, and estimated cost for three archived game tasks. Godot import and startup provide a bounded artifact check. Each condition has only one historical run, the sessions precede the current controller revision, and there is no common gameplay acceptance score. The report consequently treats the usage differences as descriptive evidence while explaining the implemented workflow that could produce them.

\begin{thebibliography}{1}
\bibitem{yoon2026pat}
Yoon et al.
\newblock PaT: Planning-after-Trial for Efficient Test-Time Code Generation.
\newblock In \emph{Proceedings of the 64th Annual Meeting of the Association for Computational Linguistics}, 2026.
\newblock \url{https://aclanthology.org/2026.acl-long.1703/}.
\end{thebibliography}

\end{document}
